\documentclass[12pt,a4paper]{article}
\usepackage[utf8]{inputenc}
\usepackage[T1]{fontenc}
\usepackage{amsmath,amssymb,amsthm}
\usepackage{mathtools}
\usepackage{geometry}
\usepackage{titlesec}
\usepackage{fancyhdr}
\usepackage{indentfirst}
\usepackage{enumitem}
\usepackage{array}
\usepackage{booktabs}
\usepackage{longtable}
\usepackage{multicol}

\geometry{margin=1in}

\pagestyle{fancy}
\fancyhf{}
\fancyhead[LE,RO]{\thepage}
\fancyhead[RE,LO]{Cayley --- Collected Mathematical Papers, Vol.~XI}
\renewcommand{\headrulewidth}{0pt}

\setlength{\parindent}{1.5em}
\setlength{\parskip}{0pt}

\newcommand{\sym}[1]{\operatorname{#1}}
\DeclareMathOperator{\sn}{sn}
\DeclareMathOperator{\cn}{cn}
\DeclareMathOperator{\dn}{dn}
\DeclareMathOperator{\am}{am}

\begin{document}

\begin{center}
{\Large THE COLLECTED MATHEMATICAL PAPERS OF}\\[6pt]
{\LARGE ARTHUR CAYLEY, Sc.D., F.R.S.,}\\[3pt]
{\large LATE SADLEIRIAN PROFESSOR OF PURE MATHEMATICS}\\
{\large IN THE UNIVERSITY OF CAMBRIDGE.}\\[12pt]
{\Large VOL.~XI.}\\[24pt]
{\large Pages 251--300}
\end{center}

\newpage


\section*{754.}
\addcontentsline{toc}{section}{754. On the Connexion of Certain Formulae in Elliptic Functions.}

\begin{center}
\textbf{ON THE CONNEXION OF CERTAIN FORMULAE IN ELLIPTIC FUNCTIONS.}
\end{center}

\begin{flushright}
\textit{[From the Messenger of Mathematics, vol.~ix. (1880), pp.~23--25.]}
\end{flushright}

\noindent
which, except as regards the determination of the constants, is the required equation
for $\log \Theta u$.

Next, differentiating twice the equation for $\Pi(u, a)$, and once the equation obtained
for $\dfrac{\Theta' u}{\Theta u}$, we have
%
\[
k^2 \sn a \cn a \dn a \, \frac{d}{du}
\left( \frac{\sn^2 u}{1 - k^2 \sn^2 u \sn^2 a} \right)
= \tfrac{1}{2} \frac{\Theta'' \Theta - \Theta'^2}{\Theta^2} (u - a)
- \tfrac{1}{2} \frac{\Theta'' \Theta - \Theta'^2}{\Theta^2} (u + a),
\]
%
and
%
\[
\frac{\Theta'' \Theta - \Theta'^2}{\Theta^2} \, u
= \frac{\Theta'' 0}{\Theta 0} - k^2 \sn^2 u,
\]
%
where, for shortness,
$\dfrac{\Theta'' \Theta - \Theta'^2}{\Theta^2} \, u$
is written to denote
$\dfrac{\Theta'' u \, \Theta u - (\Theta' u)^2}{\Theta^2 u}$,
and the like in the first equation; the right-hand side of the first equation therefore is
%
\[
- \tfrac{1}{2} k^2 \, \{ \sn^2 (u - a) - \sn^2 (u + a) \},
\]
%
or the equation becomes
%
\[
2 \sn a \cn a \dn a \, \frac{d}{du} \,
\frac{\sn^2 u}{1 - k^2 \sn^2 u \sn^2 a}
= \sn^2 (u + a) - \sn^2 (u - a),
\]
%
that is,
%
\[
\frac{4 \sn u \sn' u \sn a \cn a \dn a}
{(1 - k^2 \sn^2 u \sn^2 a)^2}
= \sn^2 (u + a) - \sn^2 (u - a).
\]

The numerator on the left-hand side must be a symmetrical function of $u$, $a$,
and hence (even if the value of $\sn' u$ were unknown) it would appear that $\sn' u$ must
be a mere constant multiple of $\cn u \dn u$; assuming, however, the actual value,
$\sn' u = \cn u \dn u$, the formula is
%
\begin{multline*}
\frac{4 \sn u \cn u \dn u \sn a \cn a \dn a}
{(1 - k^2 \sn^2 u \sn^2 a)^2} \\
\begin{aligned}
&= \sn^2 (u + a) - \sn^2 (u - a) \\
&= \{ \sn (u + a) + \sn (u - a) \} \, \{ \sn (u + a) - \sn (u - a) \}.
\end{aligned}
\end{multline*}

The factor $\{ \sn (u + a) + \sn (u - a) \}$ becomes $= 2 \sn u$ for $a = 0$,
and this suggests that the factor $\sn u$ on the left-hand side is a factor of
$\{ \sn (u + a) + \sn (u - a) \}$. That $\cn u$ is \textit{not} a factor hereof
would follow from the properties of the period $K$; viz. for $u = K$, $\cn u = 0$,
but $\{ \sn (u + a) + \sn (u - a) \}$, $= 2 \sn (K + a)$ is not $= 0$;
and, similarly, that $\dn u$ is \textit{not} a factor from the properties of the
period $iK$; hence, $\cn u$, $\dn u$ belong to the other factor
$\{ \sn (u + a) - \sn (u - a) \}$, and by symmetry $\cn a$, $\dn a$ belong to the
first-mentioned factor. And we are thus led to assume
%
\begin{align*}
\sn (u + a) + \sn (u - a) &= 2M \sn u \cn a \dn a, \\
\sn (u + a) - \sn (u - a) &= 2M' \sn a \cn u \dn u,
\end{align*}
%
where
%
\[
\text{denom.} = 1 - k^2 \sn^2 a \sn^2 u,
\]
%
and $MM' = 1$. Some further investigation is wanting to show that $M$ and $M'$
are constants, but assuming that they are so and each $= 1$, the formulae give
at once the ordinary expression for $\sn (u + a)$; that is, we have the
addition-equation for the function $\sn$.

\newpage


\section*{755.}
\addcontentsline{toc}{section}{755. On the Matrix $(a, b \\ c, d)$, and in connexion therewith the function $\frac{ax+b}{cx+d}$.}

\begin{center}
\textbf{ON THE MATRIX $\left(\begin{array}{c|c} a, & b \\ \hline c, & d \end{array}\right)$,
AND IN CONNEXION THEREWITH}\\[4pt]
\textbf{THE FUNCTION $\dfrac{ax+b}{cx+d}$.}
\end{center}

\begin{flushright}
\textit{[From the Messenger of Mathematics, vol.~ix. (1880), pp.~104--109.]}
\end{flushright}

In the preceding paper, [due to Prof.~W.~W.~Johnson,] the theory of the symbolic
powers and roots of the function $\dfrac{ax+b}{cx+d}$ is developed in a complete and satisfactory
manner; the results in the main agreeing with those obtained in the original memoir,
Babbage, ``On Trigonometrical Series,'' \textit{Memoirs of the Analytical Society} (1813), Note~I.
pp.~47--50, and which are to some extent reproduced in my ``Memoir on the Theory
of Matrices,'' \textit{Phil.\ Trans.}, t.~CXL VIII. (1858), pp.~17--37, [152]. I had recently
occasion to reconsider the question, and have obtained for the $n$th function $\phi^n x$, where
$\phi x = \dfrac{ax+b}{cx+d}$, a form which, although substantially identical with Babbage's, is a more
compact and convenient one; viz. taking $\lambda$ to be determined by the quadric equation
%
\[
\frac{(\lambda + 1)^2}{\lambda} = \frac{(a+d)^2}{ad-bc},
\]
%
the form is
%
\[
\phi^n(x) = \frac{(\lambda^{n+1} - 1)(ax+b) + (\lambda^n - \lambda)(-dx+b)}
{(\lambda^{n+1} - 1)(cx+d) + (\lambda^n - \lambda)(cx-a)}.
\]

The question is, in effect, that of the determination of the $n$th power of the
matrix $\left(\begin{array}{c|c} a, & b \\ \hline c, & d \end{array}\right)$; viz. in the notation of matrices
%
\[
(x_1, y_1) = \left(\begin{array}{c|c} a, & b \\ \hline c, & d \end{array}\right) (x, y),
\]
%
means the two equations $x_1 = ax + by$, $y_1 = cx + dy$; and then if $x_2$, $y_2$ are derived in
like manner from $x_1$, $y_1$, that is, if $x_2 = ax_1 + by_1$, $y_2 = cx_1 + dy_1$, and so on,
$x_n$, $y_n$ will be linear functions of $x$, $y$; say we have $x_n = a_n x + b_n y$,
$y_n = c_n x + d_n y$: and the $n$th power of $\left(\begin{array}{c|c} a, & b \\ \hline c, & d \end{array}\right)$
is, in fact, the matrix $\left(\begin{array}{c|c} a_n, & b_n \\ \hline c_n, & d_n \end{array}\right)$.

In particular, we have
%
\[
\left(\begin{array}{c|c} a, & b \\ \hline c, & d \end{array}\right)^{\!2}
= \left(\begin{array}{c|c} a_2, & b_2 \\ \hline c_2, & d_2 \end{array}\right)
= \left(\begin{array}{c|c} a^2+bc, & b(a+d) \\ \hline c(a+d), & d^2+bc \end{array}\right),
\]
%
and hence the identity
%
\[
\left(\begin{array}{c|c} a, & b \\ \hline c, & d \end{array}\right)^{\!2}
- (a+d) \left(\begin{array}{c|c} a, & b \\ \hline c, & d \end{array}\right)
+ (ad-bc) \left(\begin{array}{c|c} 1, & 0 \\ \hline 0, & 1 \end{array}\right) = 0;
\]
%
viz. this means that the matrix
%
\[
\left(\begin{array}{c|c}
a_2 - (a+d)a + ad-bc, & b_2 - (a+d)b \\[4pt]
c_2 - (a+d)c, & d_2 - (a+d)d + ad-bc
\end{array}\right)
= \left(\begin{array}{c|c} 0, & 0 \\ \hline 0, & 0 \end{array}\right),
\]
%
or, what is the same thing, that each term of the left-hand matrix is $=0$; which is
at once verified by substituting for $a_2$, $b_2$, $c_2$, $d_2$ their foregoing values.

The explanation just given will make the notation intelligible and show in a
general way how a matrix may be worked in like manner with a single quantity:
the theory is more fully developed in my Memoir above referred to. I proceed
with the solution in the algorithm of matrices. Writing for shortness
$M = \left(\begin{array}{c|c} a, & b \\ \hline c, & d \end{array}\right)$,
the identity is
%
\[
M^2 - (a+d)M + (ad-bc) = 0,
\]
%
the matrix $\left(\begin{array}{c|c} 1, & 0 \\ \hline 0, & 1 \end{array}\right)$ being in the theory regarded as $=1$;
viz.\ $M$ is determined by a quadric equation; and we have consequently $M^n =$
a linear function of $M$. Writing this in the form
%
\[
M^n - AM + B = 0,
\]
%
the unknown coefficients $A$, $B$ can be at once obtained in terms of $\alpha$, $\beta$, the roots
of the equation
%
\[
u^2 - (a+d)u + ad - bc = 0,\]
%
viz. we have
\begin{align*}
\alpha^n - A\alpha + B &= 0, \\
\beta^n - A\beta + B &= 0;
\end{align*}
or more simply from these equations, and the equation for $M^n$, eliminating $\alpha$, $\beta$, we
have
%
\[
\left|\begin{array}{ccc}
M^n, & M, & 1 \\
\alpha^n, & \alpha, & 1 \\
\beta^n, & \beta, & 1
\end{array}\right| = 0;
\]
%
that is,
%
\[
(\beta - \alpha) M^n - (\beta^n - \alpha^n) M + (\beta\alpha^n - \alpha\beta^n) = 0,
\]
%
or say
%
\[
M^n = \frac{\beta^n - \alpha^n}{\beta - \alpha} M
- \frac{\beta\alpha^n - \alpha\beta^n}{\beta - \alpha},
\]
%
which is the expression for $M^n$.

We may present this in a more convenient form by introducing instead of
$\alpha$, $\beta$ their differences from $\tfrac{1}{2}(a+d)$; viz. writing
%
\[
\alpha = \tfrac{1}{2}(a+d) + \theta, \qquad
\beta = \tfrac{1}{2}(a+d) - \theta,
\]
%
whence $\alpha + \beta = a + d$, $\alpha\beta = \tfrac{1}{4}(a+d)^2 - \theta^2 = ad-bc$, that is,
$\theta^2 = \tfrac{1}{4}(a-d)^2 + bc$; also $\beta - \alpha = -2\theta$, and therefore
%
\[
\frac{\beta^n - \alpha^n}{\beta - \alpha}
= \frac{(\tfrac{1}{2}(a+d) + \theta)^n - (\tfrac{1}{2}(a+d) - \theta)^n}{2\theta},
\]
%
which denoting by $P$, and writing also $Q$ for the value obtained therefrom by
changing $\theta$ into $-\theta$, that is,
%
\[
Q = \frac{(\tfrac{1}{2}(a+d) - \theta)^n - (\tfrac{1}{2}(a+d) + \theta)^n}{-2\theta},
\]
we have $P = Q$, and the formula becomes
%
\[
M^n = PM - \tfrac{1}{2} (a+d) \cdot P + \tfrac{1}{2} Q \cdot (a+d),
\]
%
that is,
%
\[
M^n = P \cdot \{M - \tfrac{1}{2}(a+d)\} + \tfrac{1}{2}(a+d) \cdot Q.
\]

But $M - \tfrac{1}{2}(a+d)$ denotes the matrix
%
\[
\left(\begin{array}{c|c}
\tfrac{1}{2}(a-d), & b \\[4pt]
c, & -\tfrac{1}{2}(a-d)
\end{array}\right),
\]
%
and we have
%
\[
M^n = P \cdot \left(\begin{array}{c|c}
\tfrac{1}{2}(a-d), & b \\[4pt]
c, & -\tfrac{1}{2}(a-d)
\end{array}\right)
+ \tfrac{1}{2}(a+d) \cdot Q \cdot
\left(\begin{array}{c|c} 1, & 0 \\ \hline 0, & 1 \end{array}\right),
\]
%
where
%
\[
P = \frac{(\tfrac{1}{2}(a+d) + \theta)^n - (\tfrac{1}{2}(a+d) - \theta)^n}{2\theta},
\qquad
Q = \frac{(\tfrac{1}{2}(a+d) + \theta)^n + (\tfrac{1}{2}(a+d) - \theta)^n}{2},
\]
%
and $\theta^2 = \tfrac{1}{4}(a-d)^2 + bc$.

The matrix equation gives
\begin{align*}
a_n &= P \cdot \tfrac{1}{2}(a-d) + \tfrac{1}{2}(a+d) \cdot Q, \\
b_n &= Pb, \\
c_n &= Pc, \\
d_n &= -P \cdot \tfrac{1}{2}(a-d) + \tfrac{1}{2}(a+d) \cdot Q,
\end{align*}
%
and we thence obtain
%
\[
\frac{a_n x + b_n}{c_n x + d_n}
\]
%
for the value of $\phi^n x$.

Writing for shortness
%
\[
\tfrac{1}{2}(a+d) = f, \qquad \tfrac{1}{2}(a-d) = g, \qquad \theta^2 = g^2 + bc,
\]
%
we have
%
\begin{multline*}
\phi^n(x) = \bigg/
\frac{(f+\theta)^n - (f-\theta)^n}{2\theta} \cdot (gx+b)
+ \frac{(f+\theta)^n + (f-\theta)^n}{2} \cdot x \\
\bigg/ \frac{(f+\theta)^n - (f-\theta)^n}{2\theta} \cdot (cx-g)
+ \frac{(f+\theta)^n + (f-\theta)^n}{2},
\end{multline*}
%
or, what is the same thing,
%
\[
\phi^n(x) = \frac{\{(f+\theta)^n - (f-\theta)^n\}(gx+b)
+ \theta\{(f+\theta)^n + (f-\theta)^n\}x}
{\{(f+\theta)^n - (f-\theta)^n\}(cx-g)
+ \theta\{(f+\theta)^n + (f-\theta)^n\}}.
\]

The case where $\theta = 0$, that is $(a-d)^2 + 4bc = 0$, requires special consideration.
Here $f = \tfrac{1}{2}(a+d)$, $g = \tfrac{1}{2}(a-d)$, $g^2 + bc = 0$, and the formulae become
%
\begin{align*}
\phi^n(x) &= \frac{nf^{n-1}(gx+b) + f^n \cdot x}{nf^{n-1}(cx-g) + f^n} \\
&= \frac{nf^{n-1}(gx+b) + f^n \cdot x}{nf^{n-1}(cx-g) + f^n}.
\end{align*}

We may write the general formula in a more convenient form by introducing
in place of $\theta$ the parameter $\lambda$, where
%
\[
\lambda = \frac{f+\theta}{f-\theta}, \qquad \text{that is} \qquad
\theta = f \cdot \frac{\lambda - 1}{\lambda + 1},
\]
%
and we thence obtain
%
\[
\phi^n(x) = \frac{(\lambda^n - 1)(gx+b) + (\lambda^{n-1} - \lambda^{-1})f \cdot x}
{(\lambda^n - 1)(cx-g) + (\lambda^{n-1} - \lambda^{-1})f},
\]
%
or substituting for $g$, $f$ their values $\tfrac{1}{2}(a-d)$, $\tfrac{1}{2}(a+d)$, this is
%
\[
\phi^n(x) = \frac{(\lambda^{n+1}-1)(ax+b) + (\lambda^n - \lambda)(-dx+b)}
{(\lambda^{n+1}-1)(cx+d) + (\lambda^n - \lambda)(cx-a)},
\]
%
which agrees with the formula originally given.

The solution may be obtained in another form by means of the quadratic equation
for the determination of $\alpha$, $\beta$; we have $\alpha$, $\beta$ the roots of
%
\[
u^2 - (a+d)u + (ad-bc) = 0,\]
%
and thence
%
\[
\alpha + \beta = a+d, \qquad \alpha\beta = ad-bc.
\]

\newpage


Writing for shortness $M = \left(\begin{array}{c|c} a, & b \\ \hline c, & d \end{array}\right)$,
and assuming $M^n = \left(\begin{array}{c|c} a_n, & b_n \\ \hline c_n, & d_n \end{array}\right)$,
we have the identity
%
\[
M^2 - (a+d)M + (ad-bc) = 0,
\]
%
that is,
%
\[
M^n = AM - B,
\]
%
where $A$ and $B$ are given as above; and thence multiplying by
$\left(\begin{array}{c|c} x, & 1 \end{array}\right)$, or say by $(x, 1)$, and observing that
%
\[
\left(\begin{array}{c|c} a_n, & b_n \\ \hline c_n, & d_n \end{array}\right)
\left(\begin{array}{c} x \\ \hline 1 \end{array}\right)
= \left(\begin{array}{c} a_n x + b_n \\ \hline c_n x + d_n \end{array}\right),
\]
%
we have
%
\[
(a_n x + b_n,\; c_n x + d_n) = A(ax+b,\; cx+d) - B(x,\; 1),
\]
%
that is,
\begin{align*}
a_n x + b_n &= A(ax+b) - Bx, \\
c_n x + d_n &= A(cx+d) - B,
\end{align*}
%
and thence
%
\[
\phi^n x = \frac{a_n x + b_n}{c_n x + d_n}
= \frac{A(ax+b) - Bx}{A(cx+d) - B}.
\]

In verification, observe that for $n=0$, $A=0$, $B=-1$, and the formula gives
$\phi^0 x = \dfrac{x}{1} = x$; for $n=1$, $A=1$, $B=0$, and the formula gives
$\phi x = \dfrac{ax+b}{cx+d}$.

Writing the formula in the form
%
\[
\phi^n x = \frac{(ax+b) - \frac{B}{A} \cdot x}{(cx+d) - \frac{B}{A}},
\]
%
and observing that $\dfrac{B}{A}$ is a symmetrical function of $\alpha$, $\beta$, we see that
$\phi^n x$ depends only on $\alpha^n$, $\beta^n$, that is on $\alpha^n + \beta^n$ and $\alpha^n \beta^n$;
but $\alpha\beta = ad-bc$, whence $(\alpha\beta)^n = (ad-bc)^n$, and we have
%
\[
\alpha^n + \beta^n = (ad-bc)^n \left(\frac{1}{\alpha^n} + \frac{1}{\beta^n}\right),
\]
%
that is, the same function of $a$, $d$ as of $d$, $a$; and $(ad-bc)^n$ is symmetrical
in regard to $a$, $d$; hence $\alpha^n + \beta^n$, and therefore $\dfrac{B}{A}$, is symmetrical in
regard to $a$, $d$; that is, we have $\phi^n(x;a,d) = \phi^n(x;d,a)$.

In particular, $n=2$ gives
%
\[
\phi^2 x = \frac{a(ax+b) + b(cx+d)}{c(ax+b) + d(cx+d)}
= \frac{(a^2+bc)x + b(a+d)}{c(a+d)x + (d^2+bc)},
\]
%
which agrees with the foregoing value.

The equation $M^n = AM - B$ may also be obtained as follows; we have
%
\[
M = \left(\begin{array}{c|c} a, & b \\ \hline c, & d \end{array}\right),
\qquad
M^n = \left(\begin{array}{c|c} a_n, & b_n \\ \hline c_n, & d_n \end{array}\right),
\]
%
and thence $M^n - AM + B = 0$ gives
\begin{align*}
a_n &= Aa - B, & b_n &= Ab, \\
c_n &= Ac, & d_n &= Ad - B,
\end{align*}
%
where $A$ and $B$ are the same linear functions of $a_n$, $d_n$ as of $a$, $d$;
viz. we have
%
\[
A = \frac{\beta^n - \alpha^n}{\beta - \alpha}, \qquad
B = \frac{\beta\alpha^n - \alpha\beta^n}{\beta - \alpha},
\]
%
and thence
%
\[
\phi^n x = \frac{A(ax+b) - Bx}{A(cx+d) - B},
\]
%
as above.

It is to be remarked that if $\lambda^m - 1 = 0$, where $m$, the least exponent for
which this equation is satisfied, is for the moment taken to be greater than $2$,
the terms in $[\;]$ are
%
\[
(\lambda-1)(ax+b) + (1-\lambda)(-dx+b),
\]
and
%
\[
(\lambda-1)(cx+d) + (1-\lambda)(cx-a);
\]
viz. these are $(\lambda-1)(a+d)x$, and $(\lambda-1)(a+d)$, or if $a+d = 0$, they
both vanish; that is, if $a+d = 0$, then $\phi^m x = x$, or the function is periodic
of the order $m$. If $a+d = 0$, then $\lambda = -1$, and the period is $=2$; we have
%
\[
\phi x = \frac{ax+b}{cx-a},
\qquad
\phi^2 x = \frac{a(ax+b) + b(cx-a)}{c(ax+b) + d(cx-a)} = x,
\]
%
since $d = -a$.

But the case $m=1$ is a very remarkable one; we have here $\lambda = 1$, and
the relation between the coefficients is thus $(a+d)^2 = 4(ad-bc)$, or what is
the same thing $(a-d)^2 + 4bc = 0$.

And then determining the values for $\lambda = 1$ of the vanishing fractions
which enter into the formulae, we find
\begin{align*}
a_n x + b_n &= \tfrac{1}{2}(a+d)^{n-1} \{(n+1)(ax+b) + (n-1)(-dx+b)\}, \\
c_n x + d_n &= \tfrac{1}{2}(a+d)^{n-1} \{(n+1)(cx+d) + (n-1)(cx-a)\},
\end{align*}
or as these may also be written
\begin{align*}
a_n x + b_n &= \tfrac{1}{2}(a+d)^{n-1} [x\{n(a-d) + (a+d)\} + 2nb], \\
c_n x + d_n &= \tfrac{1}{2}(a+d)^{n-1} [x \cdot 2nc + \{-n(a-d) + a + d\}],
\end{align*}
which for $n=0$ become as they should do $a_0 x + b_0 = x$,
$c_0 x + d_0 = 1$, and for $n=1$ they become
%
\[
a_1 x + b_1 = ax + b, \qquad c_1 x + d_1 = cx + d.
\]
%
We thus do not have $\phi^n x = x$, and the function is not periodic of
the order $1$; but we have
%
\[
\phi^n x = \frac{x\{n(a-d) + (a+d)\} + 2nb}
{x \cdot 2nc + \{-n(a-d) + a + d\}},
\]
%
which for $n = \infty$ becomes $= \dfrac{(a-d)x + 2b}{2cx - (a-d)}$, or writing herein
$(a-d)^2 + 4bc = 0$, that is $bc = -\tfrac{1}{4}(a-d)^2$, the value is
%
\[
\frac{(a-d)x + 2b}{2cx - (a-d)} = \frac{(a-d)x + 2b}{-\tfrac{1}{2}(a-d)^2 \cdot \frac{1}{b} \cdot x - (a-d)}
= \frac{(a-d)x + 2b}{-(a-d)\left\{\frac{(a-d)x}{2b} + 1\right\}},
\]
%
which is independent of $x$, and equal to either of these equal quantities;
and if from these two values of $u$ we eliminate $\lambda$, we obtain for $u$ the
quadric equation
%
\[
cu^2 - (a-d)u - b = 0,
\]
%
that is,
%
\[
u = \frac{au+b}{cu+d},\]
%
as is, in fact, obvious from the consideration that $n$ being indefinitely large
the $n$th and $(n+1)$th functions must be equal to each other. In the latter
case, as $\lambda^n$ is indefinitely small, we have the like formulae, and we obtain
for $u$ the same quadric equation: the two values of $u$ are however not the
same, but (as is easily shown) their product is $= -b \div c$; $u$ is therefore
the other root of the quadric equation. Hence, as $n$ increases, the function
$\phi^n x$ continually approximates to one or the other of the roots of this quadric
equation. The equation has equal roots if $(a-d)^2 + 4bc = 0$, which is
the relation existing in the above-mentioned special case; and here
$u = \dfrac{1}{2c}(a-d)$, $= \dfrac{-2b}{a-d}$, which result is also given by the formulae
of the special case on writing therein $n = \infty$.

\newpage


\section*{756.}
\addcontentsline{toc}{section}{756. A Geometrical Construction relating to Imaginary Quantities.}

\begin{center}
\textbf{A GEOMETRICAL CONSTRUCTION RELATING TO IMAGINARY QUANTITIES.}
\end{center}

\begin{flushright}
\textit{[From the Messenger of Mathematics, vol.~x. (1881), pp.~1--3.]}
\end{flushright}

Let $\alpha$, $\beta$, $\gamma$ be real quantities, and consider the imaginary quantity
$\alpha + \beta\omega + \gamma\omega^2$, where $\omega$ is an imaginary cube root of unity, so that
$1 + \omega + \omega^2 = 0$, and therefore $\omega^2 + \omega + 1 = 0$.

We have
%
\[
(\alpha + \beta\omega + \gamma\omega^2)
(\alpha + \beta\omega^2 + \gamma\omega)
= \alpha^2 + \beta^2 + \gamma^2 - \beta\gamma - \gamma\alpha - \alpha\beta,
\]
%
which may be written
%
\[
= \tfrac{1}{2} \{ (\beta-\gamma)^2 + (\gamma-\alpha)^2 + (\alpha-\beta)^2 \};
\]
%
and we thence have the theorem that the product of the two conjugate
imaginaries is $=0$ if $\alpha = \beta = \gamma$, but is in every other case positive.

If for a moment $\omega$ denotes generally an imaginary $n$th root of unity, and
we consider the imaginary $\alpha + \beta\omega + \gamma\omega^2 + \ldots$, then multiplying this by
the like imaginary with $\omega^{-1}$, $\omega^{-2}$, $\ldots$ in place of $\omega$, $\omega^2$, $\ldots$ respectively,
the product is
%
\[
\alpha^2 + \beta^2 + \gamma^2 + \ldots
+ 2\beta\gamma \cos(\omega^1 + \omega^{-1})
+ 2\gamma\alpha \cos(\omega^2 + \omega^{-2})
+ \ldots,
\]
%
which in the particular case $n=3$ becomes
%
\[
\alpha^2 + \beta^2 + \gamma^2 - \beta\gamma - \gamma\alpha - \alpha\beta,
\]
%
as above.

We may geometrically represent the imaginary $\alpha + \beta\omega + \gamma\omega^2$
by means of a point $P$ in a plane; viz. taking an origin $O$ and
axes $Ox$, $Oy$, if through $P$ we draw lines $PA = \alpha$, $PB = \beta$, $PC = \gamma$
inclined at angles of $120^{\circ}$ to each other, so that $PA$, $PB$, $PC$ are in fact
measured along the sides of an equilateral triangle, then the coordinates of $P$
are given by
%
\[
x = \alpha + \beta\cos 120^{\circ} + \gamma\cos 240^{\circ}
= \alpha - \tfrac{1}{2}\beta - \tfrac{1}{2}\gamma,
\]
and
%
\[
y = \beta\sin 120^{\circ} + \gamma\sin 240^{\circ}
= \tfrac{1}{2}\sqrt{3}(\beta - \gamma).
\]

The conjugate imaginary is represented in like manner by a point $P'$
obtained by reflecting $P$ in the axis of $x$; and then $OP^2 = OP'^2$ is equal
to the product of the two conjugate imaginaries, which as above is
%
\[
\tfrac{1}{2}\{ (\beta-\gamma)^2 + (\gamma-\alpha)^2 + (\alpha-\beta)^2 \}.
\]

Hence $OP^2 = 0$, that is $O$ and $P$ coincide, only in the case $\alpha = \beta = \gamma$;
and except in this case, $OP^2$ is positive.

The representation thus agrees with the fundamental theorem that the
product of an imaginary by its conjugate is in every case positive, except
only when the two factors vanish together.

In the case where $\alpha$, $\beta$, $\gamma$ are positive, and one of them, say $\gamma$, is $=0$,
the point $P$ lies on the line $OA$, and the distance $OP = \alpha - \beta$, is positive
or negative according as $\alpha > \beta$ or $\alpha < \beta$; if $\alpha = \beta$, then $P$ coincides with $O$.

\newpage


\section*{757.}
\addcontentsline{toc}{section}{757. On a Smith's Prize Question, relating to Potentials.}

\begin{center}
\textbf{ON A SMITH'S PRIZE QUESTION, RELATING TO POTENTIALS.}
\end{center}

\begin{flushright}
\textit{[From the Messenger of Mathematics, vol.~xi. (1882), pp.~15--18.]}
\end{flushright}

A \textit{spherical shell} is attracted according to the law of the inverse square
of the distance; to find the potential at any external or internal point.

The solution depends of course on the value of the solid angle which the
shell subtends at the point in question; and the determination of this solid
angle is effected by means of the spherical excess of the cone which stands
on the shell, that is of the cone whose vertex is at the point and whose base
is the rim of the shell.

Let $a$ be the radius of the sphere, $c$ the distance of the centre from the
point, $\theta$ the half-angle of the cone which the rim subtends at the centre of
the sphere, and $V$ the potential. Then if $d\Sigma$ is an element of the surface,
and $r$ its distance from the point, we have
%
\[
V = \int \frac{d\Sigma}{r},
\]
%
and the integral is to be taken over the spherical shell.

To evaluate this, let $P$ be the point, $O$ the centre of the sphere; take
$OP$ as the axis of $z$, and let the coordinates of a point on the sphere be
$(a\sin\theta\cos\phi,\; a\sin\theta\sin\phi,\; a\cos\theta)$.
Then if $d\Sigma$ is the element of surface at this point, we have
%
\[
d\Sigma = a^2 \sin\theta \, d\theta \, d\phi,
\]
%
and the distance $r$ is given by
%
\[
r^2 = a^2 + c^2 - 2ac\cos\theta.
\]

Hence for a complete cone, that is for the whole sphere, we should have
%
\[
V = \int_0^{2\pi} \!\! \int_0^{\pi}
\frac{a^2 \sin\theta \, d\theta \, d\phi}{\sqrt{a^2 + c^2 - 2ac\cos\theta}}
= 2\pi a^2 \int_0^{\pi} \frac{\sin\theta \, d\theta}{\sqrt{a^2 + c^2 - 2ac\cos\theta}}.
\]

Let $u = \cos\theta$, then $du = -\sin\theta \, d\theta$, and the integral becomes
%
\[
V = 2\pi a^2 \int_{-1}^{1} \frac{du}{\sqrt{a^2 + c^2 - 2acu}}
= \frac{2\pi a^2}{ac} \left[ -\sqrt{a^2 + c^2 - 2acu} \right]_{-1}^{1}
= \frac{2\pi a}{c} \{ -(a-c) + (a+c) \},
\]
%
that is $V = \dfrac{4\pi a^2}{c}$ for the whole sphere, which is the known result.

For a spherical shell bounded by two cones of half-angles $\alpha$ and $\beta$
($\alpha > \beta$), we have
%
\[
V = 2\pi a^2 \int_{\beta}^{\alpha}
\frac{\sin\theta \, d\theta}{\sqrt{a^2 + c^2 - 2ac\cos\theta}}
= \frac{2\pi a}{c} \{ \sqrt{a^2 + c^2 - 2ac\cos\beta}
- \sqrt{a^2 + c^2 - 2ac\cos\alpha} \}.
\]

If the point is external ($c > a$), and the shell is a zone bounded by
parallels of latitude, this gives the potential at the external point; and if
the point is internal ($c < a$), it gives the potential at the internal point.

In the case where the shell is a cap, that is where $\alpha = 0$ (the axis of
the cap), we have $\cos\alpha = 1$, and
%
\[
V = \frac{2\pi a}{c} \{ \sqrt{a^2 + c^2 - 2ac\cos\beta} - (c-a) \},
\]
if $c > a$; and
%
\[
V = \frac{2\pi a}{c} \{ (a-c) - \sqrt{a^2 + c^2 - 2ac\cos\beta} \},
\]
if $c < a$.

\newpage


\section*{758.}
\addcontentsline{toc}{section}{758. Solution of a Senate-House Problem.}

\begin{center}
\textbf{SOLUTION OF A SENATE-HOUSE PROBLEM.}
\end{center}

\begin{flushright}
\textit{[From the Messenger of Mathematics, vol.~xi. (1882), pp.~25--27.]}
\end{flushright}

Given that
%
\[
\frac{x}{a} + \frac{y}{b} + \frac{z}{c} = 0,
\qquad
\frac{a}{x} + \frac{b}{y} + \frac{c}{z} = 0,
\]
show that
%
\[
\frac{x^2}{a^2} + \frac{y^2}{b^2} + \frac{z^2}{c^2}
= \frac{y^2 z^2}{b^2 c^2} + \frac{z^2 x^2}{c^2 a^2} + \frac{x^2 y^2}{a^2 b^2}.
\]

From the given equations we have
%
\[
\frac{x}{a} + \frac{y}{b} = -\frac{z}{c},
\qquad
\frac{a}{x} + \frac{b}{y} = -\frac{c}{z},
\]
%
and multiplying these together,
%
\[
1 + \frac{bx}{ay} + \frac{ay}{bx} + 1 = 1,
\]
%
that is
%
\[
\frac{bx}{ay} + \frac{ay}{bx} + 1 = 0,
\qquad \text{or} \qquad
\frac{x^2}{a^2} + \frac{xy}{ab} + \frac{y^2}{b^2} = 0.
\]

Similarly we find
%
\[
\frac{y^2}{b^2} + \frac{yz}{bc} + \frac{z^2}{c^2} = 0,
\qquad
\frac{z^2}{c^2} + \frac{zx}{ca} + \frac{x^2}{a^2} = 0.
\]

From these three equations we may eliminate $\dfrac{yz}{bc}$,
$\dfrac{zx}{ca}$, $\dfrac{xy}{ab}$; or more simply, observe that we have
%
\[
\frac{x^2}{a^2} + \frac{y^2}{b^2} + \frac{z^2}{c^2}
+ \frac{yz}{bc} + \frac{zx}{ca} + \frac{xy}{ab} = 0.
\]

Now from the second given equation,
%
\[
\frac{ayz + bzx + cxy}{xyz} = 0,
\qquad \text{that is} \qquad
ayz + bzx + cxy = 0,
\]
%
and therefore
%
\[
\frac{yz}{bc} + \frac{zx}{ca} + \frac{xy}{ab} = 0.
\]

Hence
%
\[
\frac{x^2}{a^2} + \frac{y^2}{b^2} + \frac{z^2}{c^2} = 0.
\]

But this is not the required result; we have to show that it is equal to
%
\[
\frac{y^2 z^2}{b^2 c^2} + \frac{z^2 x^2}{c^2 a^2} + \frac{x^2 y^2}{a^2 b^2}.
\]

Now from $\dfrac{x^2}{a^2} + \dfrac{xy}{ab} + \dfrac{y^2}{b^2} = 0$, we have
%
\[
\frac{x^2 y^2 z^2}{a^2 b^2 c^2}
= \frac{z^2}{c^2} \left( -\frac{xy}{ab} \right)
= \frac{z^2}{c^2} \left( \frac{x^2}{a^2} + \frac{y^2}{b^2} \right),
\]
%
and similarly for the other terms; whence adding, the required result follows.

\newpage

\section*{759.}
\addcontentsline{toc}{section}{759. Illustration of a Theorem in the Theory of Equations.}

\begin{center}
\textbf{ILLUSTRATION OF A THEOREM IN THE THEORY OF EQUATIONS.}
\end{center}

\begin{flushright}
\textit{[From the Messenger of Mathematics, vol.~xi. (1882), pp.~34--36.]}
\end{flushright}

If in the equation
%
\[
x^n - p_1 x^{n-1} + p_2 x^{n-2} - \ldots + (-1)^n p_n = 0,
\]
we substitute for $x$ the values $\alpha$, $\beta$, $\gamma$, \ldots (the roots), and multiply
the resulting equations together, we obtain a symmetric function of the roots
which may be expressed in terms of the coefficients.

In particular, for the cubic equation
%
\[
x^3 - p_1 x^2 + p_2 x - p_3 = 0,
\]
let $\alpha$, $\beta$, $\gamma$ be the roots; then substituting $\alpha$, $\beta$, $\gamma$ successively,
and multiplying, we obtain
%
\[
(\alpha^3 - p_1 \alpha^2 + p_2 \alpha - p_3)
(\beta^3 - p_1 \beta^2 + p_2 \beta - p_3)
(\gamma^3 - p_1 \gamma^2 + p_2 \gamma - p_3) = 0.
\]

But since $\alpha^3 = p_1 \alpha^2 - p_2 \alpha + p_3$, the first factor vanishes; and similarly
for the others; so that the product is $=0$, as it should be.

More generally, if we form the equation whose roots are $\alpha^2$, $\beta^2$, $\gamma^2$,
this may be done by eliminating $x$ between the given equation and $y = x^2$;
or what is the same thing, writing the cubic in the form
%
\[
x(x^2 + p_2) = p_1 x^2 + p_3,
\]
%
and squaring, we obtain
%
\[
x^2 (x^2 + p_2)^2 = (p_1 x^2 + p_3)^2,
\]
%
that is, writing $y = x^2$,
%
\[
y(y + p_2)^2 = (p_1 y + p_3)^2,
\]
%
or
%
\[
y^3 + (2p_2 - p_1^2) y^2 + (p_2^2 - 2p_1 p_3) y - p_3^2 = 0.
\]

The coefficients of this equation are thus expressed as symmetric functions
of $\alpha$, $\beta$, $\gamma$; and we may verify that
%
\begin{align*}
\alpha^2 + \beta^2 + \gamma^2 &= p_1^2 - 2p_2, \\
\beta^2\gamma^2 + \gamma^2\alpha^2 + \alpha^2\beta^2 &= p_2^2 - 2p_1 p_3, \\
\alpha^2 \beta^2 \gamma^2 &= p_3^2.
\end{align*}

\newpage


\section*{760.}
\addcontentsline{toc}{section}{760. Reduction of $\int dx / \sqrt{1-x^4}$ to Elliptic Integrals.}

\begin{center}
\textbf{REDUCTION OF $\displaystyle\int \frac{dx}{\sqrt{1-x^4}}$ TO ELLIPTIC INTEGRALS.}
\end{center}

\begin{flushright}
\textit{[From the Messenger of Mathematics, vol.~xi. (1882), pp.~43--45.]}
\end{flushright}

The integral $\displaystyle\int \frac{dx}{\sqrt{1-x^4}}$ may be reduced to the standard
form of elliptic integral by various substitutions.

\textbf{First method.} Write $x^2 = \sin\theta$, then $2x\,dx = \cos\theta\,d\theta$, that is
$dx = \dfrac{\cos\theta\,d\theta}{2\sqrt{\sin\theta}}$, and
%
\[
\sqrt{1-x^4} = \sqrt{1-\sin^2\theta} = \cos\theta.
\]
%
Hence
%
\[
\int \frac{dx}{\sqrt{1-x^4}}
= \int \frac{d\theta}{2\sqrt{\sin\theta}}
= \frac{1}{2} \int \frac{d\theta}{\sqrt{\sin\theta}}.
\]

Now write $\theta = \tfrac{1}{2}\pi - 2\phi$, then $d\theta = -2\,d\phi$, and
%
\[
\sin\theta = \cos 2\phi = 1 - 2\sin^2\phi,
\]
%
so that
%
\[
\int \frac{dx}{\sqrt{1-x^4}}
= -\int \frac{d\phi}{\sqrt{1-2\sin^2\phi}},
\]
%
which is an elliptic integral of the first kind with modulus $k = \sqrt{2}$;
but as $k > 1$, we may write $\sin\phi = \tfrac{1}{\sqrt{2}} \sin\psi$, and thus reduce it
to the standard form with modulus $< 1$.

\textbf{Second method.} A more direct reduction is obtained by the substitution
%
\[
x = \frac{\cos\theta}{\sqrt{1+\sin^2\theta}}.
\]

We then have
%
\[
1 - x^4 = 1 - \frac{\cos^4\theta}{(1+\sin^2\theta)^2}
= \frac{(1+\sin^2\theta)^2 - \cos^4\theta}{(1+\sin^2\theta)^2}
= \frac{2\sin^2\theta(2-\sin^2\theta)}{(1+\sin^2\theta)^2},
\]
and
%
\[
dx = -\frac{\sin\theta(2+\sin^2\theta)}{(1+\sin^2\theta)^{3/2}}\,d\theta.
\]

Hence
%
\[
\frac{dx}{\sqrt{1-x^4}}
= -\frac{2+\sin^2\theta}{\sqrt{2}(1+\sin^2\theta)\sqrt{2-\sin^2\theta}}\,d\theta,
\]
%
which may be further reduced.

\textbf{Third method.} The substitution
%
\[
x = \frac{1-y}{1+y}
\]
%
gives
%
\[
1 - x^4 = \frac{(1+y)^4 - (1-y)^4}{(1+y)^4}
= \frac{8y(1+y^2)}{(1+y)^4},
\qquad
dx = \frac{-2\,dy}{(1+y)^2},
\]
%
and therefore
%
\[
\int \frac{dx}{\sqrt{1-x^4}}
= -\int \frac{(1+y)\,dy}{2\sqrt{2y(1+y^2)}}.
\]

Writing $y = t^2$, this becomes
%
\[
-\int \frac{(1+t^2)\,dt}{\sqrt{2(1+t^4)}}.
\]

The reduction of this integral to the standard form of Legendre may be
completed by the usual methods.

\newpage


\section*{761.}
\addcontentsline{toc}{section}{761. On the Theorem of the Finite Number of the Covariants of a Binary Quantic.}

\begin{center}
\textbf{ON THE THEOREM OF THE FINITE NUMBER OF THE COVARIANTS}\\[3pt]
\textbf{OF A BINARY QUANTIC.}
\end{center}

\begin{flushright}
\textit{[From the Messenger of Mathematics, vol.~xi. (1882), pp.~56--60.]}
\end{flushright}

The theorem that a binary quantic of any order has only a finite number of
irreducible covariants was first established by Gordan. The proof depends on
the representation of the covariants by means of symbolical products, and on
the theory of the representations of numbers as sums of triangular numbers.

For a binary quantic of order $n$, say
%
\[
(a, b, c, \ldots \mspace{-2mu}\between x, y)^n
= a x^n + n b x^{n-1} y + \tfrac{1}{2}n(n-1) c x^{n-2} y^2 + \ldots,
\]
%
the number of linearly independent covariants of degree $\theta$ and order $m$
may be represented in certain cases by the coefficient of $x^{\frac{1}{2}(n\theta-m)}$
in the expansion of $(1-x)^{-\theta}$, or by the number of partitions of
$\tfrac{1}{2}(n\theta-m)$ into $\theta$ or fewer parts.

The condition for the existence of such covariants is that $n\theta - m$ should
be even and non-negative, and that $m$ should not exceed $n\theta$.

The total number of linearly independent covariants of degree $\theta$ is
the coefficient of $x^{\frac{1}{2}n\theta}$ in the expansion of
%
\[
(1-x)^{-\theta} + x(1-x)^{-\theta} + x^2(1-x)^{-\theta} + \ldots,
\]
%
or more simply, it may be obtained from the generating function
%
\[
\frac{1}{(1-\alpha x^n)(1-\alpha x^{n-2})(1-\alpha x^{n-4})\cdots},
\]
%
where the expansion is to be made in powers of $\alpha$ and $x$, and the coefficient
of $\alpha^{\theta} x^m$ gives the number of covariants of degree $\theta$ and order $m$.

\textbf{Case of the Cubic.} For the cubic $(a,b,c,d \mspace{-2mu}\between x,y)^3$, the generating
function is
%
\[
\frac{1}{(1-\alpha x^3)(1-\alpha x)(1-\alpha x^{-1})}.
\]

The irreducible covariants are:
\begin{enumerate}[label=(\arabic*)]
\item The cubic itself, of degree $1$ and order $3$;
\item The Hessian, of degree $2$ and order $2$;
\item The cubicovariant, of degree $3$ and order $3$;
\item The discriminant, of degree $4$ and order $0$.
\end{enumerate}

These are connected by a syzygy of degree $6$ and order $6$; and there are
no other irreducible covariants.

\textbf{Case of the Quartic.} For the quartic $(a,b,c,d,e \mspace{-2mu}\between x,y)^4$, the
generating function is
%
\[
\frac{1}{(1-\alpha x^4)(1-\alpha x^2)(1-\alpha)(1-\alpha x^{-2})}.
\]

The irreducible covariants are:
\begin{enumerate}[label=(\arabic*)]
\item The quartic itself, of degree $1$ and order $4$;
\item The Hessian, of degree $2$ and order $4$;
\item The catalecticant (or $I$ invariant), of degree $2$ and order $0$;
\item The cubicovariant, of degree $3$ and order $6$;
\item The $J$ invariant, of degree $3$ and order $0$.
\end{enumerate}

These are connected by a syzygy; and the complete system consists of these
five irreducible forms.

\textbf{General Theory.} In the general case of a quantic of order $n$, the number
of irreducible covariants is finite. This may be shown by considering the
semi-invariants or sources of the covariants. A semi-invariant is a function
of the coefficients which is annihilated by the operator
%
\[
a \frac{\partial}{\partial b} + 2b \frac{\partial}{\partial c} + 3c \frac{\partial}{\partial d} + \ldots,
\]
%
and the source of a covariant is the coefficient of the highest power of $x$
in the covariant.

If we have a system of semi-invariants, we may form from them by the
operation of the $U$ operator (which raises the order) the corresponding
covariants; and conversely, every covariant has a semi-invariant for its
source.

The theorem of finitude then follows from the theorem that the number of
irreducible semi-invariants is finite; and this in turn may be proved by the
method of Gordan, or by the more general theory of Hilbert on the finite
basis of systems of algebraic forms.

In fact, if we consider the system of all semi-invariants of a given binary
quantic, this system has a finite basis; that is, there exists a finite number
of semi-invariants $J_1, J_2, \ldots, J_k$ such that every semi-invariant is a rational
integral function of these basis elements. And corresponding to this finite
basis of semi-invariants, we have a finite system of irreducible covariants.

The proof given by Gordan depends essentially on the theory of binary
forms and their symbolical representation; while the later proof by Hilbert
is applicable to forms in any number of variables, and establishes the more
general theorem of the finite basis for any system of algebraic forms.

\newpage


\section*{762.}
\addcontentsline{toc}{section}{762. On Schubert's Method for the Contacts of a Line with a Surface.}

\begin{center}
\textbf{ON SCHUBERT'S METHOD FOR THE CONTACTS OF A LINE WITH A SURFACE.}
\end{center}

\begin{flushright}
\textit{[From the Messenger of Mathematics, vol.~xi. (1882), pp.~71--84.]}
\end{flushright}

The theory of the contacts of a line with a surface of order $n$ has been
developed by Schubert by means of his calculus of enumerative geometry.
The principal results relate to the number of lines which have a contact
of given order with the surface; and the method employed is that of
symbolic notation and the principle of conservation of number.

\textbf{1. Contacts of the first order.} A line meets a surface of order $n$ in $n$
points. If two of these points coincide, the line has a contact of the first
order (simple contact) with the surface. The condition for this contact is
one condition; that is, the lines having a simple contact with the surface
form a system of $\infty^3$ lines in space (since the lines in space form a system
of $\infty^4$, and one condition reduces the number of parameters by $1$).

The number of such lines which pass through a given point, or which lie
in a given plane, may be determined by Schubert's methods.

\textbf{2. Contacts of the second order.} If three of the points of intersection
coincide, the contact is of the second order (stationary contact, or contact
at an inflexion). This imposes two conditions; and the lines having such
contact form a system of $\infty^2$ lines.

The surface of order $n$ has in general a certain number of parabolic points,
at which the second order contact is specially situated; and the locus of
these points is the parabolic curve, of order $4n(n-2)$, on the surface.

\textbf{3. Contacts of the third order.} If four points of intersection coincide,
the contact is of the third order (close contact). This imposes three
conditions; and the lines having such contact form a system of $\infty^1$ lines,
that is a congruence.

The number of lines of close contact which meet a given line, or which
pass through a given point, may be determined. For a surface of order $n$,
the number of close contacts through a point is in general $n(n-1)(n-2)$.

\textbf{4. Contacts of the fourth order (osculating contact).} If all five points of
intersection coincide, the contact is of the fourth order, or osculating. This
imposes four conditions; and the osculating lines form a complex of
$\infty^2$ lines.

The order of this complex, that is the number of osculating lines through
a given point, is for a surface of order $n$:
%
\[
N = n(n-1)(n-2)(n^2 + 3n - 2).
\]

\textbf{5. Five-point contact.} If the five points of intersection coincide in such a
way that the line has five consecutive points in common with the surface,
the contact is of the fifth order. This requires five conditions; and the
number of such lines is finite.

For a surface of order $n$, the number of lines having five-point contact is
%
\[
\tfrac{1}{2} n(n-1)(n-2)(n-3)(n^4 + 2n^3 + 11n^2 + 10n + 8).
\]

\textbf{6. Six-point contact.} For six-point contact, six conditions are required;
and in general there are no such lines on an arbitrary surface of order $n$.
But for special surfaces (ruled surfaces, developables, etc.) such contacts
may exist.

\textbf{Schubert's symbolic notation.} The conditions for contact are represented
symbolically as follows:
\begin{itemize}
\item $\epsilon$: the condition that a line meets a given line;
\item $\epsilon'$: the condition that a line lies in a given plane;
\item $p$: the condition that a line passes through a given point;
\item $g$: the condition that a line lies in a given plane through a given point.
\end{itemize}

The contact conditions are then denoted by symbols such as $\epsilon_2$,
$\epsilon_3$, etc., representing contacts of the second and third orders respectively;
and the principal formulas of Schubert are:
\begin{align*}
\epsilon_2 &= 2\epsilon - p - g, \\
\epsilon_3 &= 3\epsilon - 2p - 2g + \text{conditions}, \\
\epsilon_4 &= \text{combinations of lower contact conditions}.
\end{align*}

The exact formulae depend on the particular problem, and are verified by
the principle of conservation of number, according to which the number of
geometric configurations satisfying given conditions is independent of the
special position of the elements, provided the problem is properly posed.

\textbf{Application to the quadric surface.} For a quadric surface ($n=2$), the
only contacts are simple contacts; there are no stationary or higher contacts
on a general quadric. The lines having simple contact with the quadric
are the generators; and through each point of the quadric pass two such
lines.

\textbf{Application to the cubic surface.} For a cubic surface ($n=3$), there are:
\begin{itemize}
\item 27 lines on the surface (these have simple contact at every point of the line);
\item a finite number of lines having stationary contact;
\item the inflexional tangents, forming a congruence.
\end{itemize}

The 27 lines are classical; they may be classified according to their
configuration on the surface, and their contacts with the surface are of the
simplest kind.

\textbf{General formulae.} For a surface of order $n$, the number of lines having
contact of order $r$ ($r=1,2,3,4$) and satisfying the remaining conditions
to make the problem determinate, is given by formulae involving $n$ and
the Schubert conditions. The general formula for the number of lines
having simple contact and meeting four given lines is
%
\[
N_1 = \tfrac{1}{2} n(n-1)(n^2 + n + 2).
\]

For stationary contact and meeting three given lines:
%
\[
N_2 = n(n-1)(n-2)(3n^2 + 6n - 4).
\]

For close contact and meeting two given lines:
%
\[
N_3 = \tfrac{1}{2} n(n-1)(n-2)(n^3 + 3n^2 + 14n - 16).
\]

For osculating contact and meeting one given line:
%
\[
N_4 = \tfrac{1}{2} n(n-1)(n-2)(n^4 + 2n^3 + 11n^2 + 10n + 8).
\]

These formulae agree with those obtained by Schubert, and have been
verified by various independent methods.

\newpage


\section*{763.}
\addcontentsline{toc}{section}{763. On the Theorems of the 2, 4, 8, and 16 Squares.}

\begin{center}
\textbf{ON THE THEOREMS OF THE 2, 4, 8, AND 16 SQUARES.}
\end{center}

\begin{flushright}
\textit{[From the Messenger of Mathematics, vol.~xi. (1882), pp.~89--95.]}
\end{flushright}

The theorems relating to the expression of a number as a sum of 2, 4, 8,
or 16 squares have been established by Jacobi and others by means of
the theory of elliptic functions. The following investigation gives a direct
algebraical proof of these theorems, and exhibits the connexion between
them.

\textbf{Theorem of the Two Squares.} The number of representations of a number
$n$ as a sum of two squares is $4(d_1 - d_3)$, where $d_1$ is the number of divisors
of $n$ of the form $4m+1$, and $d_3$ the number of divisors of the form $4m+3$.

This may be written
%
\[
r_2(n) = 4 \sum_{d \mid n} (-1)^{\frac{1}{2}(d-1)},
\]
%
the summation extending over all odd divisors of $n$.

\textbf{Theorem of the Four Squares.} Every number is a sum of four squares;
and the number of representations is
%
\[
r_4(n) = 8 \sum_{d \mid n,\; 4 \nmid d} d,
\]
%
the summation extending over all divisors of $n$ which are not divisible
by $4$.

\textbf{Theorem of the Eight Squares.} The number of representations of $n$ as
a sum of eight squares is
%
\[
r_8(n) = 16 \sum_{d \mid n} (-1)^{n+d} d^3.
\]

\textbf{Theorem of the Sixteen Squares.} The number of representations of $n$ as
a sum of sixteen squares is
%
\[
r_{16}(n) = \tfrac{32}{17} \sum_{d \mid n} (-1)^{n+d} d^7
+ \tfrac{64}{17} \sum_{\substack{d \mid n \\ d \text{ even}}} d^7.
\]

\textbf{Proof by means of theta-functions.} The proofs of these theorems depend
on the properties of the theta-functions. We have
%
\[
\vartheta_3(0, q) = \sum_{n=-\infty}^{\infty} q^{n^2}
= 1 + 2q + 2q^4 + 2q^9 + \ldots,
\]
and therefore
%
\[
\vartheta_3^2(0, q) = \sum_{n=0}^{\infty} r_2(n) q^n,
\qquad
\vartheta_3^4(0, q) = \sum_{n=0}^{\infty} r_4(n) q^n,
\]
%
and similarly for 8 and 16 squares.

Now the theta-functions satisfy the differential equation
%
\[
4q \frac{d}{dq} \log \vartheta_3(0, q)
= \vartheta_2^4(0, q) + \vartheta_4^4(0, q),
\]
%
and by means of this and the known relations between the theta-functions,
the values of $r_2(n)$, $r_4(n)$, $r_8(n)$, $r_{16}(n)$ may be determined.

\textbf{Direct algebraic proof.} Let
%
\[
\sigma_k(n) = \sum_{d \mid n} d^k
\]
%
denote the sum of the $k$th powers of the divisors of $n$; and let
%
\[
\sigma_k^*(n) = \sum_{d \mid n,\; 2 \nmid d} d^k
= \sum_{d \mid n} d^k - \sum_{d \mid n,\; 2 \mid d} d^k.
\]

Then the theorems may be stated as follows:
\begin{align*}
r_2(n) &= 4\{\sigma_0^*(n; 4m+1) - \sigma_0^*(n; 4m+3)\}, \\
r_4(n) &= 8\sigma_1^*(n) - 32\sigma_1^*(\tfrac{1}{4}n), \\
r_8(n) &= 16(-1)^n \sigma_3^*(n) + 256\sigma_3^*(\tfrac{1}{4}n), \\
r_{16}(n) &= \tfrac{32}{17} \{(-1)^n \sigma_7^*(n) + 2\sigma_7^*(\tfrac{1}{2}n)
+ 16\sigma_7^*(\tfrac{1}{4}n)\},
\end{align*}
where $\sigma_k^*(\tfrac{1}{4}n)$ is understood to be $0$ when $n$ is not divisible by $4$.

The proofs depend essentially on the identities between the theta-functions,
and on the expansion of powers of $\vartheta_3(0, q)$ in series of powers of $q$.

\textbf{Connexion with modular functions.} The theorems of the squares are also
connected with the theory of modular functions. If we write
%
\[
q = e^{\pi i \tau}, \qquad \tau = \omega + i\omega',
\]
%
then $\vartheta_3(0, q)$, $\vartheta_2(0, q)$, $\vartheta_4(0, q)$ are modular functions of $\tau$; and the
identities between them are equivalent to relations between modular forms.

In particular, the theorem of four squares is equivalent to the statement
that the theta-function $\vartheta_3^4(0, q)$, which is a modular form of weight $2$, has
a development in which the coefficient of $q^n$ is determined by the divisor
function.

\textbf{Application to the representation of primes.} A prime $p$ of the form
$4m+1$ can be expressed as a sum of two squares in essentially one way;
and the number of representations is $8$ (counting order and signs).

A prime of the form $4m+3$ cannot be expressed as a sum of two squares.
Every prime can be expressed as a sum of four squares; and the number
of representations is $8(p+1)$ for an odd prime $p$.

\textbf{Special cases.} For $n=1$, we have $r_2(1) = 4$, corresponding to
$(\pm 1)^2 + 0^2$ and $0^2 + (\pm 1)^2$; $r_4(1) = 24$; $r_8(1) = 16$; $r_{16}(1) = 32$.

For $n=2$, we have $r_2(2) = 4$, from $(\pm 1)^2 + (\pm 1)^2$;
$r_4(2) = 24$; $r_8(2) = 16$; $r_{16}(2) = 32$.

These values agree with the general formulae; and the theorems have
been verified in numerous particular cases.

\vspace{1cm}
\begin{center}
\rule{0.3\textwidth}{0.4pt}
\end{center}

\end{document}
